% 放在参考文献之后。
% 所需宏包：booktabs、array。
% 整体替换旧版附录，不要追加到旧版后面。
% 每个语言版本的主稿只使用一次 \appendix。

\clearpage
\appendix

% 附录A
\section{复现与测量}
\label{app:protocol}

\paragraph{复现材料。} 可复现性声明所述 GitHub 发布包将提供可执行源码、
配置与依赖清单、运行命令、图表生成脚本和机器可读结果表，以便重新生成
论文中的结果。缺失或不可评价字段将明确报告，不以插值替代。

\paragraph{数据来源与命名。} 本文使用公开的
\emph{HumDial-FDBench}测试归档；正式论文为Wang等人
\citep{humdial2026}，数据页单独引用为\citep{humdialdataset2026}。
官方划分为Train 8,898、Dev 1,800和Test 5,000条，覆盖九类场景；公开测试归档
另含4,100个用于评分的clean文件。本文冻结的Hugging Face revision为
\texttt{ed1770a6265ffe9f0e5cd999290936ccab43482b}，数据卡标注Apache-2.0。
冻结归档中的WAV为单声道、16-bit PCM、16\,kHz；本文直接使用发布的WAV字节和
JSON标注，不做项目内重采样、降噪、响度归一化或合成混音。样本ID保留归档路径
\texttt{test/\{en,cn\}\_test\_nondev/\{scenario\}/\{base\_id\}[\_add].wav}，
clean文件是独立的评分参照。运行时cut从选定完整WAV在实验固定的非空状态边界生成，
属于项目派生回放切面，不是官方HumDial-FDBench片段或榜单分数；cut ID、时间戳、
来源成员和哈希绑定在运行manifest中。该发布物是人工录制的表演语音，不是TTS；
公开论文和数据卡没有说明去标识化或说话人匿名化流程，因此本文不宣称音频已经去标识化。

全文统一使用：\emph{FD-Bench}（Peng等）是独立基准，不使用无连字符简称；
“HumDial-FDBench”仅作为HumDial官方数据集标识保留；
\emph{Full-Duplex-Bench}指Lin等人的家族，其中v1和v1.5是离线版本，
\emph{Full-Duplex-Bench-v2}是后续实时框架；\emph{HumDial-FDBench}是基于v1.5协议
构建的HumDial基准，既不是FD-Bench数据集，也不是Full-Duplex-Bench-v2。
本文改编的事件定义统一称为“FD-Bench风格”。


冷恢复实验使用相同的Lychee-FD模型与本地Token2Wav后端
\citep{lychee2026}，
在两张A100-SXM4-40GB GPU上分别部署模型执行和一个声学lane。
共享权重保持常驻，
每组比较内部匹配模型、采样、数值、输入和播放配置。
复现需要实际执行的源码、依赖与配置清单，包括归档补丁，
不能仅依据仓库commit。

最终矩阵包含12个来源组、四种策略和两次重复，共96次执行。
E1包含44次，E2包含52次，两个环境分别统计。
测量批次包含24个待处理语音token，以及20,032个24\,kHz缓存PCM样本。
这些数值描述预定义非空切面，而非自然会话状态分布。
名义挂起时间为30秒，但时延采用实际事件计量。
新音频渲染从Worklet恢复开始，
计量至同一AudioContext中首个新计算PCM帧被渲染，
缓存音频续播不属于新计算。主指标为
\[
T_{\mathrm{first\text{-}new\text{-}PCM}}
=t_{\mathrm{new}}-t_{\mathrm{W,resume}}。
\]
Save使用安全暂停到资源就绪的区间，Ready使用RESUME接受到执行发布的区间，
二者均采用服务端单调时钟；Gap使用浏览器时钟定义为
$\max(0,t_{\mathrm{new}}-t_{\mathrm{buffered\text{-}end}})$。
这些端点不是串行片段，不能相加。

配置统计先在同来源、同环境内汇总具有对应端点的重复，
再取来源间中位数。配对差先匹配来源、重复和环境；缺失值不以零替代。
冷恢复、传输、吞吐量和交互批次分别保留身份，不合并统计。

最终矩阵包含96条计划运行，其中92条通过技术合同。Rebuild保留1条
E1暂停后原生浏览器帧重叠取消；Hybrid保留3条E1失败（两条帧重叠
取消和一条冻结切面/首个PCM观测轮次不匹配）。Hot和Joint没有失败，
E2全部通过技术合同。Hot和Rebuild没有对应的检查点、pinned缓冲或
GPU工作区测量事件，主表分别标记为N/A或NM，不填零。system pinned、
$L_{\mathrm{useful}}$和$J(T)$不属于本轮冷恢复合同，明确标记为未评价。


% 附录B
\section{联合状态与恢复验证}
\label{app:state_contract}

表\ref{tab:app_state_contract}总结联合恢复边界。
计算状态可以恢复，而已接受输入、取消决定与交付记录保持当前值。
这一区分在请求粒度上应用一致快照与回滚恢复原则
\citep{chandy1985,elnozahy2002}。

\begin{table}[tbp]
    \centering
    \small
    \setlength{\tabcolsep}{5pt}
    \caption{
        状态覆盖与恢复义务。
        外部记录持续保留，不被旧检查点副本覆盖。
    }
    \label{tab:app_state_contract}
    \begin{tabular}{
        @{}
        >{\raggedright\arraybackslash}p{0.20\linewidth}
        >{\raggedright\arraybackslash}p{0.73\linewidth}
        @{}
    }
        \toprule
        状态类别 & 内容与恢复义务 \\
        \midrule
        模型
        & 有效KV、非KV历史、处理器／控制状态、采样上下文和逻辑位置。
          将数值恢复至合法映射，并重建活动绑定。 \\
        桥接
        & 已产生但未消费的token、部分缓冲、
          初始化／flush状态和待提交输出。
          保留内容与消费顺序。 \\
        声学
        & StreamingDecoder与Token2Wav连续状态、
          Flow / HiFT缓存、实际随机上下文和待提交PCM。
          恢复兼容的执行上下文。 \\
        身份与进度
        & 请求、generation、状态版本、cut、lease epoch、sequence、
          producer epoch和各阶段水位。
          根据当前生命周期验证归属。 \\
        外部记录
        & 已接受输入、取消决定及交付／渲染区间。
          保留当前事实，拒绝过期提交。 \\
        \bottomrule
    \end{tabular}
\end{table}

捕获阶段先建立一致切面与独立主机检查点，
再交还请求私有资源。
恢复阶段准备合法映射、重建绑定并验证各参与阶段状态，
之后才允许执行和输出。
发布时重新检查当前generation、版本和取消状态；
失败或已取消的恢复不被发布。
恢复已保存的输入消费位置，不会覆盖检查点之后收到的输入。

请求身份回归区分物理输出顺序$[A,B]$
与仅包含$B$的采样选择。
request ID查找取得$B$对应的side output，
而不是物理第零行的输出。
独立归档检查还覆盖桥接待处理内容、发布前拒绝和请求局部取消。
完整token与波形比较以匹配的输入／控制历史及数值配置为条件，
不代表语义正确，也不保证任意kernel的无条件确定性。

\paragraph{协议证明义务。}
协议信封$q$包含五个字段：\textit{request\_id}、\textit{generation}、\textit{state\_version}、
\textit{cut\_id}和\textit{lease\_epoch}；每条记录再附带\textit{sequence}和
\textit{producer\_epoch}。证明采用如下归纳。初始状态只有一个归属者，
已产生集和已消费集均为空。SAVE在切面线性化点前关闭全部写入者，因此检查点内记录具有相同$q$；
排空在途提交后，检查点与队列恰好划分已产生未消费集合。RESUME只有在读取当前账本水位后才安装
新的lease epoch。之后每个提交只有在完整信封等于当前permit时接受，账本更新与队列位置推进原子完成。

因此，请求ID不匹配的记录会被拒绝，归属保持不变；token要么在检查点中，要么由消费者事务恰好移除一次，
桥接守恒保持；输入、取消和交付水位只增不减，账本单调性保持。generation、state version与lease epoch
共同拒绝迟到worker及ABA lease复用。消费者只接受下一个sequence，故PCM前缀／顺序由归纳得到。
PUBLISH是离开RESTORING的唯一转换，且只在这些检查完成后开放，因此无效状态不能执行或可见。
任一前提失败都转入\textsc{Aborted}，不会发布半恢复状态，也不会回滚外部事实。

两个线性化点在trace中可观测：\textsc{Cut}是原子关闭并快照的事件，\textsc{Restore}是原子安装新lease
和permit的事件。取消在线性化时写入账本；恢复若观察到该记录则中止，而更晚的取消会在下一次提交前使
permit失效。上述规则分别覆盖切面不一致、版本过期、ABA lease、迟到worker、桥接重复消费、PCM乱序、
取消／恢复竞争及检查点后的新输入，不依赖物理worker身份。

\begin{algorithm}[tbp]
\caption{逻辑请求$r$的保存／恢复协议。}
\label{alg:save_resume}
\begin{algorithmic}[1]
\Require 当前信封$q$、状态$(M,B,A)$和外部账本$L$
\Ensure 发布permit，或进入不可执行的中止状态
\State 原子关闭输入、token生产和输出发布入口
\State 递增\textit{state\_version}并分配\textit{cut\_id}
\State 在支持的边界暂停$M,B,A$并排空在途提交
\If{存在信封不同于$q$的记录}
    \State 拒绝该记录并中止捕获
\EndIf
\State 在\textsc{Cut}线性化点执行$K\gets\Call{Snapshot}{M,B,A,L,q}$
\If{$\neg\Call{CheckJoint}{K}$}
    \State 保留原资源；\Return \textsc{Aborted}
\EndIf
\State 释放私有资源并将$r$标为挂起
\State 读取当前$L$；若取消或低于输入／取消／交付水位则拒绝
\State 分配新的\textit{lease\_epoch}并合并切面后的输入
\State 从$K$重绑定$M,B,A$
\If{$\neg\Call{CheckJoint}{M,B,A,L,q}$}
    \State 丢弃未发布状态；\Return \textsc{Aborted}
\EndIf
\State 在\textsc{Restore}点原子安装新lease和permit
\While{$r$仍在执行}
    \If{提交信封不等于$q$，或PCM序列不是下一个前缀值}
        \State 丢弃该提交
    \Else
        \State 原子追加账本记录并推进消费位置
    \EndIf
\EndWhile
\State \Return \textsc{Published}
\end{algorithmic}
\end{algorithm}


% 附录C
\section{补充性能结果}
\label{app:performance}

表\ref{tab:app_latency}补充正文的测量端点。
表中同时报告保存、就绪、首个新PCM渲染和间隔，
每项均保留自己的起止锚点。
各指标保留自己的起止端点，
不与其他区间相加构造新的端到端指标。

\begin{table}[tbp]
    \centering
    \small
    \setlength{\tabcolsep}{6pt}
    \caption{
        补充冷恢复测量，单位秒。
        数值为对应环境内端点的来源级中位数。N/A表示事件不适用，NM表示
        没有对应测量；二者都不是零。第一组给出最新报告中完整96次运行
        分母下的策略级有效率；失败运行仍保留在分母中。
    }
    \label{tab:app_latency}
    \begin{tabular}{@{}llcrrrr@{}}
        \toprule
        环境 & 策略 & 有效/总数 & 保存 & 就绪 & 首个新PCM渲染 & 间隔 \\
        \midrule
        \multicolumn{7}{@{}l}{\emph{所有环境（最新96次运行批次）}} \\
        All & Hot     & 96/96 & \na & \na & \na & \na \\
        All & Rebuild & 92/96 & \na & \na & \na & \na \\
        All & Joint   & 96/96 & \na & \na & \na & \na \\
        All & Hybrid  & 96/96 & \na & \na & \na & \na \\
        \midrule
        E1 & Hot     & 11/11 & 0.004 & 0.001  & 1.184 & 0.349 \\
        E1 & Rebuild & 10/11 & 1.847 & 13.317 & 14.451 & 13.616 \\
        E1 & Joint   & 11/11 & 1.387 & 0.699  & 1.899 & 1.064 \\
        E1 & Hybrid  & 8/11  & 0.655 & 0.130  & 1.349 & 0.515 \\
        \midrule
        E2 & Hot     & 13/13 & 0.004 & 0.001  & 1.168 & 0.333 \\
        E2 & Rebuild & 13/13 & 1.926 & 13.250 & 14.400 & 13.565 \\
        E2 & Joint   & 13/13 & 1.387 & 0.700  & 1.872 & 1.037 \\
        E2 & Hybrid  & 13/13 & 0.696 & 0.136  & 1.315 & 0.480 \\
        \bottomrule
    \end{tabular}
\end{table}

最终批次的首个新PCM渲染区间比缓存后供给间隔多
$20{,}032/24{,}000=0.834667$秒。这是预定义缓存播放时长，
不是额外的时延项；二者是同一交付时间线的不同视角，不是独立改善。
E2中，Joint与Hybrid的独立检查点载荷分别为
672.207和22.416\,MiB，
受管理pinned缓冲分别为256和128\,MiB。
两者均保留模型设备工作区，使模型GPU净分配增加60.489\,MiB。
检查点大小、pinned存储、进程RSS和设备分配属于不同测量，
不能直接相加。

独立传输实验包含六个来源组、12对Sync/Optimized运行。
每对的相对降幅为
$(T_i^{\mathrm{sync}}-T_i^{\mathrm{opt}})/T_i^{\mathrm{sync}}$。
其中传输到接收端就绪区间的相对降幅中位数为16.00\%，绝对减少量中位数为0.533秒，
12对均为优化路径更快。该批次没有测量Worklet恢复、首个新PCM渲染、播放结束或端到端恢复时延。
额外暂停RSS与采样峰值RSS的配对增加中位数
分别为243.2和182.2\,MiB。
历史吞吐矩阵则对各负载／并发单元的逐次配对比值取均值。
其全部配置和重复应完整保留；
新吞吐测量作为独立批次处理。


% 附录D
\section{交互与外部协议}
\label{app:interaction}

自然交互先导采用FD-Bench风格的适配定义\citep{fdbench2025}。
SRR、SIR与SRIR的分母分别保留符合条件的非打断问句、
实际打断机会和成功打断事件。
EIR记录助手过早进入用户话轮。
FSED从用户语音结束计量至匹配回复开始播放，
IRD从用户插话开始计量至助手上一段语音结束。
时延保留正负符号。
显式CPU挂起不计作语音打断，
拼接PCM也不能替代真实渲染时间线。

导出的事件列按如下方式计算
$\mathrm{SRR}=k_{\mathrm{SR}}/n_{\mathrm{SR}}$、
$\mathrm{SIR}=k_{\mathrm{SI}}/n_{\mathrm{SI}}$、
$\mathrm{SRIR}=k_{\mathrm{SRI}}/n_{\mathrm{SRI}}$和
$\mathrm{EIR}=k_{\mathrm{EI}}/n_{\mathrm{EI}}$；某指标列缺失时，该事件
不进入该指标分母。每条受控路径有10个有效问句、3个实际打断机会，
6次重复中无效／取消运行数为0。唯一非零的失败相关字段是Omni的
1次reply timeout，该记录不产生FSED观测，因此该路径FSED有效$n=9$。
正文区间是source-level bootstrap区间（对6个source group重采样
200,000次，seed为20260926）；有效bootstrap次数保留在
\texttt{interaction\_metrics\_bootstrap.csv}中。

共享Lychee/DuplexPilot在线路径属于同一实现，
不是两个独立基线。
该路径采用400\,ms块末端输入，
Omni客户端采用200\,ms先发送、再等待的输入方式；
模型和实际打断暴露状态也不同。
因此，先导描述应用行为，不用于隔离冷恢复影响。
此前听音反馈只适用于已经审阅的案例，
不扩展到未审阅先导输出。

两例vLLM-Omni原生热重附着\citep{omnidoc}
的请求至ACK分别为1.512和1.421\,ms，
首PCM接收分别为725.412和776.220\,ms。
两例记录中的PCM序号没有缺失或重复，
但首个返回PCM的产生阶段未知。
LiveServe风格预加载组件\citep{liveserve2026}
提前约9.9秒准备20和15\,MiB KV。
其就绪分别提前约56.65和100.01\,ms，
但从共同源窗口计量的首新PCM接收分别晚23.7和36.4\,ms。
这些观察保留各自原生协议，
不构成统一的外部冷恢复比较。

\begin{table*}[t]
\centering
\scriptsize
\setlength{\tabcolsep}{3pt}
\caption{受控Lychee/DuplexPilot路径和原生Omni参照路径的逐source重复结果。比例为分子/分母；FSED和IRD为带符号播放层中位数（ms），括号内为有效$n$。每次重复均未出现无效或取消运行；反向通道排除是有意的非指标事件，不是失败。}
\label{tab:interaction_repeats}
\begin{tabular}{llccccrrcc}
\toprule
配置 & source group & SRR & SIR & SRIR & EIR & FSED (ms) & IRD (ms) & timeout & 无效/取消 \\
\midrule
Lychee & 0001\_0001 & 1/1 & -- & -- & 0/1 & 3666.4 (1) & -- & 0 & 0 \\
Lychee & 0001\_0003 & 1/1 & -- & -- & 0/1 & 3762.9 (1) & -- & 0 & 0 \\
Lychee & 0001\_0015 & 1/1 & 1/1 & 1/1 & 1/2 & 5582.8 (2) & 266.6 (1) & 0 & 0 \\
Lychee & 0001\_0025 & 1/1 & 1/1 & 1/1 & 1/2 & 5575.1 (2) & 278.2 (1) & 0 & 0 \\
Lychee & 0001\_0029 & 2/2 & -- & -- & 1/2 & 5808.0 (2) & -- & 0 & 0 \\
Lychee & 0001\_0031 & 1/1 & 1/1 & 1/1 & 1/2 & 3615.5 (2) & 898.9 (1) & 0 & 0 \\
\midrule
Omni & 0001\_0001 & 1/1 & -- & -- & 1/1 & -185.9 (1) & -- & 0 & 0 \\
Omni & 0001\_0003 & 1/1 & -- & -- & 1/1 & 1324.7 (1) & -- & 0 & 0 \\
Omni & 0001\_0015 & 1/2 & -- & -- & 2/2 & 302.9 (1) & -- & 1 & 0 \\
Omni & 0001\_0025 & 1/1 & 1/1 & 1/1 & 1/2 & 1507.4 (2) & 3444.6 (1) & 0 & 0 \\
Omni & 0001\_0029 & 1/1 & 1/1 & 1/1 & 2/2 & 1351.6 (2) & 564.1 (1) & 0 & 0 \\
Omni & 0001\_0031 & 1/1 & 1/1 & 1/1 & 2/2 & 1544.5 (2) & 2068.3 (1) & 0 & 0 \\
\bottomrule
\end{tabular}
\end{table*}

